\scanpage{153}
garantie que $f^0$ soit connue quand $f$ l'est [Ex : si on prend $\tau^{\prime}=\tau^{\unsure{'|2|\prime\prime}}=0$, $2\,1_k=0$, alors $f$ doit être \uline{constante}]. C'est pourquoi je préfère mettre $f^0$ dans la structure de l'application --- i.e. l'application est un\deleted{\unsure{e}} \deleted{paire}couple
\[
 \begin{aligned}
 (f,f^0):\qquad &f:A_1^{\unsure{1}}\times A_2^{\unsure{2|1}}\longrightarrow B\\
 &f^0:A_1^0\times A_2^0\longrightarrow B^0
 \end{aligned}
\]

On voit tout de suite qu'il y a un accouplement \uline{universel}, dont le module des translations est
\[
 A_1^0\otimes A_2^0
\]
et \deleted{\unsure{qu'on}} on notera $A_1\ast A_2$, espace affine sous $A_1^0\otimes A_2^{\unsure{0|2}}$, avec $(\sigma_1,\sigma_2)$-accouplement
\[
 (x_1,x_2)\longmapsto x_1\ast x_2:\quad A_1\times A_2\longrightarrow A_1\ast A_2.
\]

On peut le réaliser comme sous espace affine de
\[
 E_1\otimes E_2\supset A_1^0\otimes A_2^0
\]
($E_i$ étant le module enveloppant de l'espace affine $A_i$), l'application universelle étant induite par
% archive sheet number: 152

